\scanpage{46}\mspage{4 - 42 -}
\hfill n$^{\circ}$ 392

morphe à un idéal de $A$. Supposons donc $n\geq 2$ ; en vertu du lemme 8, il existe un sous-module libre $L_{\unsure{o|0}}$ de $M$, de rang 1, tel que $M/L_o$ soit sans torsion ; $M/L_o$ est donc de rang $n-1$. En vertu de l'hypothèse de récurrence, il y a donc un sous-module libre $L_1$ de $M/L_o$ tel que $(M/L_o)/L_1$ soit isomorphe à un idéal de $A$. Soit $L$ l'image réciproque de $L_1$ dans $M$ ; $L/L_o$ est isomorphe à $L_1$, et puisque $L_1$ est libre, $L$ est isomorphe à $L_o\oplus L_1$, (\uline{Alg}., chap. II, 3e éd., \S~1, n$^{\circ}$ 11, prop. 21), donc libre ; comme $M/L$ est isomorphe à $(M/L_o)/L_1$, le théorème est démontré.

\uline{Remarque}.- Si $M$ est réflexif, il n'en est pas nécessairement de même de $M/L$, (exerc. \qquad).

10. \uline{Modules sur les anneaux de Dedekind}.

On suppose maintenant que $A$ soit un \uline{anneau de Dedekind} ; on sait alors que les idéaux $\underline{p}\in P$ sont \uline{maximaux} et que ce sont les seuls idéaux premiers $\ne 0$ de $A$, (\S~2, n$^{\circ}$ 2, th. 1) ; le groupe $D(A)$ s'identifie au groupe $I(A)$ des idéaux fractionnaires $\ne 0$ de $A$.

PROPOSITION 21.- \uline{Soit} $A$ \uline{un anneau de Dedekind}. \uline{Tout} $A$-\uline{module pseudo-nul est nul}. \uline{Tout homomorphisme pseudo-injectif}, (resp. \uline{pseudo-surjectif}, \uline{pseudo-bijectif}, \uline{pseudo-nul}), \uline{de} $A$-\uline{modules est injectif}, (resp. \uline{surjectif}, \uline{bijectif}, \uline{nul}).

La première assertion a déjà été démontrée, (n$^{\circ}$ 4, \uline{Exemple} 1) ; les autres s'en déduisent immédiatement.

PROPOSITION 22.- \uline{Soient} $A$ \uline{un anneau de Dedekind}, $M$ \uline{un} $A$-\uline{module de type fini}. \uline{Les propriétés suivantes sont équivalentes} :

a) $M$ \uline{est sans torsion}.

b) $M$ \uline{est réflexif}.

c) $M$ \uline{est projectif}.

On sait déjà, (sans hypothèse sur l'anneau intègre $A$), que b) implique a), (n$^{\circ}$ 2, \uline{Remarque} 1), et que c) implique b), (\uline{Alg}., chap. II, 3e éd., \S~2, n$^{\circ}$ 7,
% archive sheet number: 45
